\section{Preliminary Domain Knowledge}

During our preliminary research, we were especially interested in machine learning projects that make use of data already collected by everyday devices. Smartphones stood out as a strong example because they continuously record information from built-in sensors such as accelerometers and gyroscopes. This suggests that phones may support useful security or health-related applications without requiring extra hardware.

% when a different walk is detected, lock the phone, or trigger an alert. This is a form of behavioral biometrics, where the system learns to recognize the unique patterns of an individual's behavior—in this case, their gait.

For our project, we want to apply this idea to a practical security problem: determining whether a phone is still in its owner's control while the person carrying it is walking. This matters because theft or unauthorized access does not always happen in an obvious way. Someone might take a phone and walk away with it, but it could also be picked up briefly for a nearby action such as making a payment and then returned. In either case, detecting only the exact moment the phone is taken is usually not enough.

Our main idea is that a person's gait may act as a behavioral fingerprint. A phone carried while walking records repeated motion patterns, and we believe these signals may contain user-specific information such as stride timing, acceleration structure, and frequency patterns. We found a previous study and repository that approached a related problem as a multiclass classification task, which suggests that different people may produce distinguishable motion signatures with high predictive accuracy.

Based on this, our project aims to model walking sequences from smartphone accelerometer and gyroscope data and map them into an embedding space. We then plan to use cosine similarity or another comparison method to estimate whether a new walking segment is likely to come from the phone's owner or from someone else. In other words, the machine learning task is owner-versus-non-owner verification based on gait-related sensor data collected during walking.